% !Mode:: "TeX:UTF-8"   %%winedt 以utf8编码打开
% Can only be compiled with xelatex+bibtex

%双行标题
\documentclass[twoside,longtitle]{LZUthesis}
%单行标题
%\documentclass[twoside]{LZUthesis}

\usepackage{float}
\usepackage{fancybox}
\usepackage{calc}
\usepackage{mathdots}
\usepackage{graphicx}
\usepackage{listings}

%声明图片后缀名
\DeclareGraphicsExtensions{.pdf,.eps}

\makeatletter

% 设置图形文件的搜索路径
\graphicspath{{figures/}}

% 小节标题靠左对齐
\CTEXsetup[format+={\flushleft}]{section}

% 取消链接的颜色（黑白打印时启用）
%\hypersetup{colorlinks=false}

%使文档居中，打印时应注释掉
\evensidemargin 0.93 true cm
\setlength{\hoffset}{-0.3 cm}

%允许equationarray分页换行
\allowdisplaybreaks

%使字体清晰，且透明位图不会使页面文字粗细不一
\usepackage{eso-pic}
\AddToShipoutPicture{%
\special{pdf: put @thispage <</Group << /S /Transparency /I true /CS /DeviceRGB>> >>}%
}

%页面背景色
%\definecolor{yellow}{rgb}{0.99,0.99,0.70}
%\pagecolor{yellow}

\makeatother


\begin{document}

%分类号；一般不要求学生填写
\classification{}

%密级；申请保密后需填写
\confidential{}

%中文标题
\title{兰州大学研究生学位论文}

%中文标题第二行，题目较短时删除之，并去除文档类选项 longtitle
\titleadd{\LaTeX{}模板简介}

%英文标题
\englishtitle{Introduction to the \LaTeX{} Template of}

%英文题目第二行，题目较短时删除之，并去除文档类选项 longtitle
\englishtitleadd{Lanzhou University Thesis}

%作者汉语姓名
\author{刘海滨}

%专业
\major{物理学•理论物理}

%学位
\degree{博士}

%研究方向
\direction{XXXX}

%导师姓名+职称
\advisor{某某某~教授}

%论文工作起止年月
\datebeginAndEnd{2012年9月至2015年5月}

%论文提交日期
\submitdate{}

%答辩日期
\defenddate{}

%学位授予日期
\degreedate{}

%左侧页眉
\lzuthesis{兰州大学博士学位论文}

%生成封面
\maketitle

%生成声明页
\makestatement

%前文-罗马页码
\frontmatter\pagenumbering{Roman}

%中文摘要
\begin{abstract}
本文是兰州大学学位论文的\LaTeX{} (包含LyX)模板及其使用说明。除了介绍文档类LZUthesis的用法外，本文还是一份简要的学位论文写作指南。
\end{abstract}

%中文关键词
\keywords{兰州大学，学位论文，\LaTeX{}模板，LyX模板}

%英文摘要
\begin{englishabstract}
This paper is a \LaTeX{} (include LyX) thesis template of Lanzhou University. Besides the usage of the LaTeX document class LZUthesis, a brief guideline for writing the thesis is also included.
\end{englishabstract}

%英文关键词
\englishkeywords{Lanzhou University, Thesis, \LaTeX{} Template, LyX Template}

%自动生成章节目录以及pdf书签
\tableofcontents{}


%正文部分-数字页码
\mainmatter

%正文页眉页脚样式
\pagestyle{lzu}


\chapter{引言\label{chap:intro}}

虽然兰州大学研究生院制定了\href{http://ge.lzu.edu.cn/degree/xwsq/lwgf/201103/1186.htm}{学位论文撰写要求}(见附录~\ref{chap:standard})，但只提供了\href{http://ge.lzu.edu.cn/degree/xwsq/lwgf/201103/1187.htm}{~Word 模板}。广大理科生更加青睐的\LaTeX{}却没有标准模板。\href{http://blog.sciencenet.cn/home.php?mod=space&uid=117412&do=blog&id=512804}{LZUthesis~} 宏包是赵振华以中科院学位论文宏包\href{http://www.ctex.org/PackageCASthesis}{~CASthesis~}作为基础模板，根据兰州大学研究生院学位论文撰写要求修改的。作者对LZUthesis宏包进行了进一步修改(主要变更见附录~\ref{chap:changelog})，使之更加符合兰大学位论文的撰写要求，并制作了相应的LyX模板。

宏包底层使用的是Xe\LaTeX{}，而不是CJK方案，原因在于：
\begin{enumerate}
\item Xe\LaTeX{}支持系统字体，对中文友好。
\item Xe\LaTeX{}支持unicode，可编译各国文字。
\item Xe\LaTeX{}支持常见的图片格式，如：pdf、eps、jpg、png 等。
\end{enumerate}


模板使用LyX的优点是所见即所得，比较直观，缺点是\TeX{}中一些复杂功能难以实现。另外LyX可以导出标准的\TeX{}代码(参见第~\ref{sec:export}~ 节)，方便转入\TeX{}编辑器继续撰写。

青睐\TeX{}模板的用户，直接编辑模板根目录下的Thesis.tex即可。注意它不是Thesis.lyx直接导出的，如果要使用此模板，不要让LyX导出\TeX{}时覆盖此文件。


\chapter{系统配置}

\texttt{LZUthesis}宏包使用了宏包amsmath、amsthm、amsfonts、amssymb、bm、ctex 和hyperref。目前大多数的\TeX{}系统中都包含有这些宏包。但是中文字体在各个系统中却不尽相同：Windows 下使用C\TeX{} 套装一般不需要额外配置字体，但Linux 和Mac一般情况下需要设定中文字体，中文字体名称可通过命令~\lstinline!sudo fc-list :lang=zh-cn!~ 查看，具体配置见本章内容。请注意字体设定可能要根据你自己系统的中文字体做相应调整。


\section{Windows配置}
\begin{itemize}
\item 安装C\TeX{}套装，最新的版本可以从\href{http://www.ctex.org}{~www.ctex.org} 网站下载。
\item 安装完成后，把所有宏包升级到最新(开始菜单--CTEX--Update(Admin))。
\item 安装LyX，最新的LyX安装包可以从\href{http://www.lyx.org/Download}{~www.lyx.org/Download} 下载。
\end{itemize}

\section{Linux配置}

软件安装以ubuntu为例，其他发行版应根据具体情况安装。
\begin{itemize}
\item 安装texlive以及texlive-lang-cjk。


\doublebox{\begin{minipage}[t]{0.86\textwidth}%
\begin{lstlisting}[language=bash,breaklines=true]
sudo apt-get install texlive texlive-xetex
sudo apt-get install texlive-lang-cjk
\end{lstlisting}
%
\end{minipage}}

\item 安装LyX。


\doublebox{\begin{minipage}[t]{0.86\textwidth}%
\begin{lstlisting}[language=bash,breaklines=true]
sudo apt-get install lyx
\end{lstlisting}
%
\end{minipage}}

\item 中文字体设置

\begin{enumerate}
\item 文档类选项(文档--首选项--文档类)中添加选项{[}nofonts{]}
\item 导言区(文档--首选项--\LaTeX{}导言区)中添加


\tiny
\begin{center}
\hrule


\begin{verbatim}


\setCJKmainfont[BoldFont=文泉驿微米黑,ItalicFont=文泉驿正黑] {文泉驿正黑} \setCJKsansfont{文泉驿微米黑}
\setCJKmonofont{文泉驿等宽正黑}
\setCJKfamilyfont{zhsong}{文泉驿正黑}
\setCJKfamilyfont{zhhei}{文泉驿微米黑}
\setCJKfamilyfont{zhfs}{文泉驿等宽微米黑}
\setCJKfamilyfont{zhkai}{文泉驿等宽正黑}
\newcommand*{\songti}{\CJKfamily{zhsong}} % 宋体
\newcommand*{\heiti}{\CJKfamily{zhhei}}   % 黑体
\newcommand*{\kaishu}{\CJKfamily{zhkai}}  % 楷书
\newcommand*{\fangsong}{\CJKfamily{zhfs}} % 仿宋


\end{verbatim}


\hrule
\end{center}
\normalsize

\end{enumerate}
\end{itemize}

\section{Mac配置}
\begin{itemize}
\item 安装Mac\TeX{}，最新安装文件可以从\href{https://tug.org/mactex/}{~https://tug.org/mactex/}
网站下载。
\item 安装LyX，最新的LyX安装包可以从\href{http://www.lyx.org/Download}{~www.lyx.org/Download} 下载。
\item 中文字体设置

\begin{enumerate}
\item 文档类选项(文档--首选项--文档类)中添加选项{[}nofonts{]}
\item 导言区(文档--首选项--\LaTeX{}导言区)中添加


\tiny
\begin{center}
\hrule


\begin{verbatim}


\setCJKmainfont[BoldFont=Heiti SC Light,ItalicFont=Kaiti SC]   {Songti SC Light}
\setCJKsansfont{Heiti SC Light} \setCJKmonofont{Songti SC Light}
\setCJKfamilyfont{zhsong}{Songti SC Light}
\setCJKfamilyfont{zhhei}{Heiti SC Light}
\setCJKfamilyfont{zhfs}{STFangsong}
\setCJKfamilyfont{zhkai}{Kaiti SC}
\newcommand*{\songti}{\CJKfamily{zhsong}} % 宋体
\newcommand*{\heiti}{\CJKfamily{zhhei}}   % 黑体
\newcommand*{\kaishu}{\CJKfamily{zhkai}}  % 楷书
\newcommand*{\fangsong}{\CJKfamily{zhfs}} % 仿宋


\end{verbatim}


\hrule
\end{center}
\normalsize

\end{enumerate}
\end{itemize}


本节内容未经实测，请实测者予以确认或提出修改意见。


\chapter{模版使用}


\section{模板文件结构\label{sec:files}}

整个模板根目录的文件列表如下：

\begin{tabular}{|l|l|l|l|}
\hline
LZUthesis.cls & ---LZUthesis宏包 & \textcolor{red}{{*}} & \textcolor{magenta}{\$}\tabularnewline
\hline
LZUthesis.cfg & ---LZUthesis宏包配置文件 & \textcolor{red}{{*}} & \textcolor{magenta}{\$}\\
\hline
LZUthesis.layout & ---LZUthesis的LyX布局文件 & \textcolor{red}{{*}} & \textcolor{magenta}{\$}\\
\hline
CASmthesis.cls & ---CASthesis宏包 & \textcolor{red}{{*}} & \textcolor{magenta}{\$}\\
\hline
CASmthesis.cfg & ---CASthesis宏包配置文件 & \textcolor{red}{{*}} & \textcolor{magenta}{\$}\\
\hline
lzubib.bst & ---引文样式文件 & \textcolor{red}{{*}} & \textcolor{magenta}{\$}\\
\hline
bib/thesis.bib & ---bib数据库 & \textcolor{red}{{*}} & \textcolor{magenta}{\$}\\
\hline
figures/lzu.pdf & ---兰州大学校名标准字 & \textcolor{red}{{*}} & \textcolor{magenta}{\$}\\
\hline
Thesis.lyx & ---LyX模板主文件 & \textcolor{red}{{*}} & -\\
\hline
Envelope.lyx & ---A3封面LyX模板 & \textcolor{red}{{*}} & -\\
\hline
spine.lyx & ---书脊LyX模板 & \textcolor{red}{{*}} & -\\
\hline
Thesis.tex & ---\TeX{}模板主文件 & - & \textcolor{magenta}{\$}\\
\hline
Makefile & ---Makefile文件 & - & -\\
\hline
gbt7714-2005.bst & ---符合国标gbt7714-2005的引文样式文件 & - & -\\
\hline
\end{tabular}

注：\textcolor{red}{{*}}表示LyX模板必须的文件，\textcolor{magenta}{\$} 表示\LaTeX{}模板必须的文件。


\section{基本设置}
\begin{enumerate}
\item 模版默认使用的文档类选项(文档--首选项--文档类--自定义)是twoside,~longtitle，即双面，长标题。

\begin{enumerate}
\item 双面：此时页码、页边距以及章开始页等会对奇偶页区别设置，可更改为oneside。
\item 长标题：若标题较短时应去除选项longtitle，并删除中英文标题的第二行。
\end{enumerate}
\item 图片搜索路径默认设置为模板根目录下的figures/。
\item bib数据库默认设置为模板根目录下的bib/thesis.bib。
\end{enumerate}

\section{示例}

论文中最常使用的一些功能在本节中给出示例，其他LyX功能的使用详见帮助-- 用户手册，\TeX{}使用则可参考The \TeX{}book\cite{Knut04-T}。


\subsection{公式}

\begin{equation}
\hat{H}=\frac{\epsilon}{2}\hat{\sigma}_{z}-\frac{\Delta}{2}\hat{\sigma}_{x}+\sum_{k}\omega_{k}\hat{b}_{k}^{\dagger}\hat{b}_{k}+\sum_{k}\frac{g_{k}}{2}\hat{\sigma}_{z}(\hat{b}_{k}+\hat{b}_{k}^{\dagger})\label{eq:sbm}
\end{equation}



\subsection{表格}

\begin{table}[H]
\begin{centering}
\begin{tabular}{|c|c|c|c|c|}
\hline
Alpha & Beta & Gamma & Delta & Theta\\
\hline
$\alpha$ & $\beta$ & $\gamma$ & $\delta$ & $\theta$\\
\hline
$A$ & $B$ & $\Gamma$ & $\Delta$ & $\Theta$\\
\hline
\end{tabular}
\par\end{centering}

\protect\caption{希腊字母表\label{tab:Greek}}


\end{table}



\subsection{图形}

\begin{figure}[H]
\begin{centering}
\includegraphics[width=0.618\textwidth]{figures/lzu}
\par\end{centering}

\protect\caption{兰州大学校名标准字\label{fig:lzu}}
\end{figure}



\subsection{引用}


\subsubsection{交叉引用}

对所有需要引用的公式、表格、图形，执行插入--标签后，即可使用插入-- 交叉引用自动产生引用。
\begin{itemize}
\item 哈密顿量见方程~(\ref{eq:sbm})。
\item 希腊字母表见表~\ref{tab:Greek}。
\item 校名标准字如图~\ref{fig:lzu}。
\end{itemize}

\subsubsection{文献引用}

将引文的bib数据库(默认文件名为thesis.bib)放入模板根目录下的bib 文件夹，即可通过插入--文献引用自动产生引文。
\begin{itemize}
\item Journal：An article\cite{Liu13PRA-A}。
\item Book：C\TeX{} FAQ\cite{Wu05-C}。
\end{itemize}

\subsection{makefile}

makefile仅限\LaTeX{}模板编译使用，适合Linux和Mac用户使用。
\begin{enumerate}
\item make：只编译源文件，快速调试用。
\item make ref：包含引用的完整编译。
\item make clean：清除临时文件。
\end{enumerate}

\section{注意事项}
\begin{enumerate}
\item 模板目录下的必须文件(见第~\ref{sec:files}~节)注意不要删除。
\item 请尽量升级\TeX{}宏包到最新，一些旧包会使排版稍微不正常。截止到\today ，此模板在最新宏包下工作正常。
\item Bibkey中不能有中文字符。
\item LyX中必须用空格时，尽量用\textasciitilde{}或(插入--格式-- 强制间距~) 代替，不用空格的原因是：LyX不会中文断句，使用空格将使得LyX编辑界面混乱。
\item Hyperref已经在模板中声明过，如果在LyX中执行插入--超链接会造成冲突(直接编辑\TeX{}代码无此问题)，这时只能通过插入\TeX{} 代码的方式使用，比如：\lstinline!\href{https://www.google.com}{Google}!。
\item LyX中，enumerate环境结束后，必须先添加一个分割线环境才能换行，具体参见第\ref{chap:intro}章中列表结束后的蓝色分割线(蓝色分割线仅现于LyX)。
\end{enumerate}

\section{论文打印}


\subsection{文档设置}
\begin{enumerate}
\item 按需选取twoside或者oneside的选项。
\item 启用导言区(文档--首选项--\LaTeX{}导言区)的~\lstinline!\hypersetup{colorlinks=false}!以取消链接的颜色。
\item 预留左侧装订线：注释掉导言区中如下两行


\doublebox{\begin{minipage}[t]{0.86\textwidth}%
\begin{lstlisting}[language=bash,breaklines=true]
\evensidemargin 0.93 true cm
\setlength{\hoffset}{-0.3 cm}
\end{lstlisting}
%
\end{minipage}}

\end{enumerate}

\subsection{书脊生成}
\begin{enumerate}
\item spine.lyx生成竖排的书脊文字。
\item Envelope.lyx生成含书脊的封皮，纸张大小A3+。
\item 若中文标题中含有英文，需要将导言区(文档--首选项--\LaTeX{} 导言区)中的\lstinline!\shuji! 替换为如下样式：\footnotesize
\begin{lstlisting}
\shuji[某\hspace{0.3em}\raisebox{-5pt}{English}\hspace{-0.1em}~某]
\end{lstlisting}
\normalsize 本文标题中的\LaTeX{}即是如此处理，对于全中文标题的情况可做上面的相反替换。
\end{enumerate}

\chapter{进阶功能}


\section{导出\protect\TeX{}\label{sec:export}}

文件-导出-\LaTeX{}(Xe\TeX{})，可导出为tex文件进行进一步编辑，该tex 文件的编译顺序为: Xe\LaTeX{}--Bib\TeX{}--Xe\LaTeX{}--Xe\LaTeX{}。 注意LyX 默认会自动将图片转换为pdf格式，造成一些图片有两个版本，不处理也不会有影响。若想处理，建议保留pdf版本的，因为Xe\LaTeX{}处理pdf图片时，编译速度快；若选择保留其他格式，应在tex文件正文前声明图片的后缀名，比如：\lstinline!\DeclareGraphicsExtensions{.pdf,.eps}!


\section{导入Bibitem到正文}

如果需要在没有bib文件的情况下也可编译，则首先把导出的tex文件完整编译一次，再将生成的bbl文件内容拷贝到tex正文中替换掉如下两行即可

\doublebox{\begin{minipage}[t]{0.86\textwidth}%
\begin{lstlisting}[language=bash,breaklines=true]
\bibliographystyle{lzubib}
\bibliography{bib/thesis}
\end{lstlisting}
%
\end{minipage}}


\section[修复stemV问题]{修复stemV问题\footnote{此问题已在修复透明位图导致文字粗细不一的问题时同时得到解决。}}

由于中文字体某些参数设置不当，\LaTeX{}生成的pdf在Adobe阅读器下中文字体(特别是宋体与仿宋)会稍显模糊，可通过修改stemV值改善。
\begin{enumerate}
\item 下载并安装\href{https://www.pdflabs.com/tools/pdftk-the-pdf-toolkit/}
{~pdftk free}
\item cmd中切换到pdf所在目录，执行


\doublebox{\begin{minipage}[t]{0.86\textwidth}%
\begin{lstlisting}[language=bash,breaklines=true]
pdftk in.pdf output out.pdf uncompress
\end{lstlisting}
%
\end{minipage}}

\item 用文本编辑器编辑out.pdf，查找并修改simsun与fangsong的stemV 值为50，保存。
\item 执行以下命令，得到final.pdf


\doublebox{\begin{minipage}[t]{0.86\textwidth}%
\begin{lstlisting}[language=bash,breaklines=true]
pdftk out.pdf output final.pdf compress
\end{lstlisting}
%
\end{minipage}}

\end{enumerate}

\section{自动生成成果页}
\begin{enumerate}
\item 此方法仅适用于\LaTeX{}模板。
\item 导言中添加：


\doublebox{\begin{minipage}[t]{0.86\textwidth}%
\begin{lstlisting}[language=bash,breaklines=true]
%%自动生成发表文章目录
\usepackage[resetlabels]{multibib}
%把参考文献加入目录中，但排除目录本身
\usepackage[nottoc]{tocbibind}
\newcites{phd}{在学期间的研究成果}
\end{lstlisting}
%
\end{minipage}}

\item 在参考文献列表后添加


\doublebox{\begin{minipage}[t]{0.86\textwidth}%
\begin{lstlisting}[language=bash,breaklines=true]
%%自动生成发表文章目录
%此处插入所有发表文章的bibkey
\nocitephd{Liu13PRA-A}
%定义引文样式为lzubib
\bibliographystylephd{lzubib}
%发表文章的bib库
\bibliographyphd{bib/thesis}
\end{lstlisting}
%
\end{minipage}}

\item 此种情况使用makefile不能完整编译，但WinEdt可以完整编译。执行顺序为:


Xe\LaTeX{}--Bib\TeX{}--Xe\LaTeX{}--Xe\LaTeX{}。


命令行下手动编译时需要额外执行\lstinline!BibTeX phd.aux!。

\end{enumerate}



%附录部分
\appendix

%自动生成参考文献列表，需要模板根目录下有：bib/thesis.bib
\bibliographystyle{lzubib}
\bibliography{bib/thesis}


%发表文章目录
\begin{publications}{99}
\item Ji-Hong Zhang, \textbf{Ling-Yun Wu}, and Xiang-Sun Zhang, \textit{Reconstruction of DNA sequencing by hybridization}, \href{http://link.aps.org/doi/10.1103/PhysRevA.87.052139}{Bioinformatics, \textbf{19}(1): 14-21.} (2013).
\item \textbf{吴凌云}, \href{http://www.ctex.org/CTeXFAQ}{\textit{CTEX FAQ (常见问题集)}}, (2004).
\end{publications}


%致谢
\begin{thanks}
Ajouter de la vie aux jours lorsqu'on ne peut plus ajouter de jours
à la vie.

\begin{flushright}
---Jean Bernard
\par\end{flushright}
\end{thanks}


\chapter{主要变更\label{chap:changelog}}
\begin{enumerate}
\item 修复目录中摘要和abstract超链接的错误
\item 修复长标题在摘要部分只显示第一行
\item 修复中英文摘要第一行不缩进
\item 修复页眉上不显示完整标题
\item 页眉改为双线
\item 替换lzu.eps为清晰版本
\item 添加lyx模板
\item 修正行间距
\item 设置Twoside时页码在装订线右侧，oneside时居中
\item 调整封面、声明页布局
\item 修正章标题为三号字
\item 改变图表的编号格式为2-1
\item 修正致谢页的字体及行距
\item 修正一些中文字段，如摘要、Key words、在学期间的研究成果等
\item 修改小节标题为不缩进
\item 添加A3封面，书脊
\item 修改bst样式文件
\item 调小装订线距离为0.6cm
\item 致谢、成果偶数页留白
\item 单独定义makestatement
\item 合并cls文件，longtitle作为类选项进入
\item 更改各处全标题的调用
\item 成果页格式调整为与参考文献相同
\item 摘要重定义为chapter，间距微调
\item chapter标题上下间距微调
\item 去除摘要、Abstract、目录的页眉
\item lzu.eps替代为lzu.pdf
\item 修复透明位图导致文字粗细不一的问题
\end{enumerate}


\chapter[论文写作规范]{兰州大学研究生学位论文写作规范\label{chap:standard}}

研究生学位论文是研究生科学研究工作的全面总结，是描述其研究成果、代表其研究水平的重要学术文献资料，是申请和授予相应学位的基本依据。为规范我校研究生学位论文编写格式，根据《科学技术报告、学位论文和学术论文的编写格式》（GB/T~7713-1987）和《学位论文编写规则》（GB/T~7713.1-2006) 等国家标准，制定本研究生学位论文写作规范。


\section{学位论文基本要求}

硕士学位论文，要求对所研究的课题有新见解或新成果，并对本学科发展或经济建设、社会进步有一定意义，表明作者掌握坚实的基础理论和系统的学科知识，具有从事学术研究或担负专门技术工作的能力。学位论文应在导师指导下，由硕士研究生本人独立完成。

博士学位论文，要求对所研究的课题在材料、角度、观点、方法、理论等方面或某方面有创新性成果，并对学术发展、经济建设和社会进步有较重要的意义，表明作者掌握坚实宽广的基础理论和系统深入的学科知识，具有独立从事学术研究的能力。学位论文应在导师指导下，由博士研究生本人独立完成。

学位论文应当用规范汉字进行撰写，除古汉语研究中涉及的古文字和参考文献中引用的外文文献之外，均采用简体中文撰写。

学位论文必须是一篇系统完整的、有创造性的学术论文。

博士学位论文一般在5万字左右，硕士学位论文一般在3万字左右。


\section{学位论文基本结构}

学位论文基本结构包括前置部分、主体部分和结尾部分。


\subsection{前置部分}
\begin{enumerate}
\item 封面（包括扉页）
\item 独创性声明
\item 关于学位论文使用授权的申明
\item 摘要页（包括中英文摘要）
\item 目录页
\item 插图和附表清单（可根据需要）
\end{enumerate}

\subsection{主体部分}
\begin{enumerate}
\item 引言（或绪论）
\item 正文
\item 结论
\end{enumerate}

\subsection{结尾部分}
\begin{enumerate}
\item 参考文献
\item 附录（可根据需要）
\item 插图和附表（可根据需要）
\item 索引（可根据需要）
\item 作者简历及在学期间所取得的科研成果（可根据需要）
\item 封底
\end{enumerate}

\section{编写规范与要求}


\subsection{前置部分}
\begin{enumerate}
\item 封面：封面包括分类号、密级、校名、学位论文中文题目、英文题目、作者姓名、导师姓名、学科和专业名称、提交时间等内容（附件1）。


分类号：按中国图书分类法，根据学位论文的研究内容确定。


密级：仅限于涉密学位论文（论文课题来源于国防军工项目）填写，密级应根据涉密学位论文确定，分绝密、机密和秘密三级，并注明保密期限（保密论文的密级和保密期限由学校保密办公室审批确定）。非涉密学位论文不得填写密级。


论文题目：论文题目要准确反映整篇论文的核心内容。中文题目一般不超过25 个汉字，英文题目翻译简短准确，一般不超过150个字母，必要时可以加副标题。


学科和专业名称：按研究生培养的学科专业目录，规范填写。


学位级别：博士研究生填写“博士”，硕士研究生填写“硕士”。

\item 独创性声明：见附件2。
\item 关于学位论文使用授权的申明：见附件2。
\item 摘要：包括中文摘要和英文摘要两部份。摘要是论文内容的总结概括，应简要说明论文的研究目的、基本研究内容、研究方法、创新性成果及其理论与实际意义，突出论文的创新之处。不宜使用公式、图表，不标注引用文献。硕士论文摘要一般应在500 字左右，博士论文摘要一般应在1000 字左右。英文摘要与中文摘要内容要一致。关键词应体现论文特色，具有语义性，一般为3—5 个。（见附件3）
\item 目录页：论文中内容标题的集合。目录页中每行均由标题名称和页码组成。包括摘要、章节或大标题的序号和名称、参考文献、注释、索引等。（见附件4）。
\item 插图和附表清单：论文中如图表较多，可以分别列出清单置于目录页之后。图的清单应有序号、图题和页码。表的清单应有序号、表题和页码。
\end{enumerate}

\subsection{主体部分}

包括引言（或绪论）、正文和结论。主体部分从另页右页开始，每一章另起页。
\begin{enumerate}
\item 基本要求

\begin{enumerate}
\item 引言（绪论）：包括论文的研究目的、流程和方法等；还包括论文研究领域的历史回顾，文献回溯，理论分析等内容。
\item 正文：正文是学位论文的主体和核心部分，包括研究背景、立论根据、研究内容、研究方法与过程、研究结果与分析、研究结论及其意义。要求论述正确、逻辑严密、层次分明、文字流畅简练、公式图表清晰规范、数据真实可靠，公式推导和计算结果正确无误，避免使用文学性质的带感情色彩的非学术性词语。论文中如出现非通用性的新名词、新术语、新概念，应作相应解释。


图：图应具有“自明性”，应鲜明清晰。图包括曲线图、构造图、示意图、框图、流程图、记录图、地图、照片等。照片上应有表示目的物尺寸的标度。图的编号和图题规范，并应置于图下方。


表：表应具有“自明性”。表的编号和表题应规范，并置于表上方。表题应简单明了。


表的编排建议采用国际通行的三线表。一般是内容和测试项目由左至右横读，数据依序竖读。如某个表需要转页接排，在随后的各页上应重复表的编号。编号后跟表题（可省略）和“（续）”，置于表上方。续表均应重复表头和有关表述。


公式：论文中的公式应另行起，并缩格书写，与周围文字留足够的空间区分开。如有两个以上的公式，应用从“1”开始的阿拉伯数字进行编号，并将编号置于括号内。公式的编号右端对齐，公式与编号之间可用“…”连接。公式较多时，应分章编号。较长的公式需要转行时，应尽可能在“＝”处回行，或者在“+”、“－”、“×”等记号处回行。


公式下面的“式中：”空两个字起排，单独占一行。公式中所要解释的符号按先左后右，先上后下顺序分行空两个字排，再用破折号与释文连接，回行时与上一行释文对齐。上下行的破折号对齐。


引文标注：论文中引用的文献的标注方法可采用顺序编码制，也可采用著者－出版年制，但全文必须统一。如：


德国学者N.克罗斯研究了瑞十巴塞尔市附近侏罗山中老第三纪断裂对第三系褶皱的控制{[}25{]};之后，他又描述了西里西亚第3条大型的近南北向构造带，并提出地槽是在不均一的块体的基底上发展的思想{[}26,27{]}。(顺序编码制)


结构分析的子结构法最早是为解决飞机结构这类大型和复杂结构的有限元分析问题而发展起来的（Przemienicki,1968）(著者-出版年制)


注释：当论文中的字、词或短语，需要进一步加以说明，而又没有具体的文献来源时，用注释。注释一般在社会科学研究中用得较多。应控制论文中的注释数量，不宜过多。注释采用文中编号加“脚注”的方式，置于当页的页脚。

\end{enumerate}
\item 章节图表标号规则

\begin{enumerate}
\item 章节标号


论文章节按序分层。层次以少为宜，根据实际需要选择。各层次标题可以用阿拉伯数字，也可以用汉字。


用阿拉伯数字标题：各层次要用阿拉伯数字连续标号；不同层次的数字之间用小圆点“．”相隔，末位数字后面不加点号，如，“1.1”，“1.1.1”等；章、节编号全部顶格排，编号与标题之间空1个字符的间隙。正文另起行，前空2 个汉字符起排，回行时顶格排。例如：


第一章 ××××（章大标题）


×××××××××××××××××××××××××××


1.1 ××××（一级节标题）


×××××××××××××××××××××××××××


1.1.1 ××××（二级节标题）


1.1.1.1 ××××（根据需要，也可设三级节标题）


第二章 ××××（章大标题）


2.1 ××××（一级节标题）


2.1.1 ××××（二级节标题）


用汉字标题：各层次要按“章”、“节”、“条”、“款”等依次编号，如：“第一章”，“第一节”、“一”、“（一）”等；章、节编号居中，编号与标题之间空1 个汉字符的间隙。正文另起行，前空2个汉字符起排，回行时顶格排。例如：


第一章 ××××（章大标题）


第一节 ××××××××（节大标题）


一、×××××××××（一级节标题）


（一）×××××××××××××（二级节标题）


（二）×××××××××××××（二级节标题）


1.×××××××××××××××（三级节标题）


（1）××××××××××××

\item 图、表等标号


论文中的图、表、附注、公式、算式等，一律用阿拉伯数字分章依序连续编码。其标注形式应便于互相区别，如：图l.1（第1章第一个图）、图2.2 （第二章第二个图）；表3.2 （第三章第二个表）等。


图题置于图序之后，图序和图题之间空1个汉字符的间隙，并应置于图下方的居中位置。引用图应在图题右上角标出文献来源。


表题置于表序之后，表序和表题之间空1个汉字符的间隙，并置于表上方的居中位置。

\item 页码、页眉编写规则


学位论文页码的摘要页和目录页用罗马数字单独编连续码，正文和后置部分用阿拉伯数字编连续码。单面复印时页码排在页脚居中位置，双面复印时页码分别按左右侧排列。


页眉、页脚文字均采用小五号宋体，左侧页眉为“兰州大学博（硕）士学位论文”，右侧为学位论文题目名称；页眉下横线可为单横线也可用上粗下细文武线。

\end{enumerate}
\item 结论


论文的结论是最终的、总体的结论，不是正文中各段的小结的简单重复。结论应包括论文的核心观点，交代研究工作的局限，提出未来工作的意见或建议。结论应该准确、完整、明确、精练。


如果不能导出一定的结论，也可以没有结论而进行必要的讨论。

\end{enumerate}

\subsection{结尾部分}
\begin{enumerate}
\item 参考文献


参考文献是文中引用的有具体文字来源的文献集合，学位论文中引用他人成果之处均应如实、详细地列出参考文献目录。其著录项目和著录格式遵照《文后参考文献著录规则》（GB/T~7714－2006）的规定执行。


参考文献可以列在各章的结尾，也可以列在正文的末尾，并另起页。所有被引用文献均要列入参考文献表中。引文采用顺序编码标注时，参考文献表按编码顺序排列，引文采用著作－出版年制标注时，参考文献表应按著者字顺和出版年排序。


各种主要参考文献按如下格式编排（附件5）：


学术期刊：序号作者 文题 刊名 年 卷号（期号） 起止页码


专（译）著：序号作者（译者） 书名 出版地 出版者 出版年 起止页码


学位论文：序号作者 文题 {[}XX学位论文{]} 授予单位所在地 授予单位 授予年份 起止页码


专利：序号申请者 专利名 国名 专利文献种类 专利号 出版日期


技术标准：序号发布单位 技术标准代号 技术标准名称 出版地：出版者，出版日期


电子文献：序号 作者 出版年 题名 出版地 出版者 ［引用日期］ 获取和访问路径

\item 附录


附录作为主体部分的补充，并不是必须的。下列内容可以作为附录编于论文后：


为了整篇论文材料的完整，但编入正文又有损于编排的条理性和逻辑性，这一材料包括比正文更为详尽的信息、研究方法和技术更深入的叙述，对了解正文内容有用的补充信息等；


由于篇幅过大或取材于复制品而不便于编入正文的材料；


不便于编入正文的罕见珍贵资料；


对一般读者并非必要阅读，但对本专业同行有参考价值的资料；


某些重要的原始数据、数学推导、结构图、统计表、计算机打印输出件等。

\item 索引：根据需要可以编排分类索引，关键词索引等。
\item 作者简历：包括教育经历、工作经历、攻读学位期间发表的论文和完成的工作等。攻读学位期间发表的论文按发表的时间顺序，列出本人在攻读学位期间发表或已录用的学术论文清单，包括发表文章与出版论著（已发表或已录用）、已获专利、科研成果及获奖情况分项列出。


发表文章与出版论著列出格式同参考文献格式。


获奖成果列出格式为：获奖成果名称 奖励级别与获奖等级 获奖年份与日期 获奖人


获专利成果列出格式为：获专利名称 专利号 何国何类专利 授权日期 专利人

\item 致谢：致谢对象限于对课题研究、学位论文完成等方面有较重要帮助的人员。主要包括：对国家科学基金、资助研究工作的奖学金基金、合同单位、资助或支持的企业、组织或个人；对协助完成研究工作和提供便利条件的组织或个人；对在研究工作中提出建议和提供帮助的人；对给予转载和引用权的资料、图片、文献、研究思想和设想的所有者；对其他应感谢的组织和个人。
\end{enumerate}

\section{论文打印规格}

学位论文用A4标准纸(210~mm~×297~mm)打印、印刷或复印，按论文顺序装订成册，论文顺序依次为：封面（包括扉页）、论文原创性声明、关于学位论文使用授权的申明、中文摘要、英文摘要、目录、论文正文、参考文献、符号说明、注释、附录、攻读学位期间发表的学术论文目录、致谢等。

论文页边距一般要求：上边距3cm、下边距2.54cm，左右边距3.17cm，页眉页脚2.0cm。


\subsection{封面}

论文封面颜色：博士学位论文为浅红色，硕士学位论文（除硕士专业学位论文）为浅蓝色，硕士专业学位论文为银白色。

论文题目用三号字，宋体，加粗，其他信息用小三号字，宋体，加粗，居中。题目可严格按照封面模板格式控制各部分的字体、字号。

书脊须写明论文题目、姓名、兰州大学、年份，用四号字，仿宋，行距16 磅，段前段后0磅。


\subsection{摘要}
\begin{enumerate}
\item 中文摘要

\begin{enumerate}
\item 论文题目为三号字，黑体，加粗，居中，单倍行距，段前24 磅，段后18 磅。可分行，每行左右两边至少留3个汉字符空格。
\item 论文题目下为“摘要”，两字间空1个汉字符间隙，三号字，黑体，加粗，居中，单倍行距，段前24磅，段后18磅。
\item 摘要内容为小四号字，宋体，行距20磅，段前段后0磅，首行缩进2 个汉字符。
\item 摘要内容后下空1行，行距20磅，段前段后0磅。
\item 摘要内容后为“关键词”，小四号字，宋体，行距20磅，段前段后0 磅，“关键词”三字加粗。关键词数量为3-5个，词与词间用逗号隔开。
\end{enumerate}
\item 英文摘要

\begin{enumerate}
\item 除“Abstract”外，论文中的英文一律采用“Times~New~Roman” 字体。论文英文题目全部采用大写字母，16磅，加粗，居中，单倍行距，段前24 磅，段后18磅。论文题目可分行，每行左右两边至少留6 个字符空格。
\item 英文摘要为“Abstract”，Arial字体，16磅，加粗，居中，单倍行距，段前24磅，段后18磅。
\item 英文摘要内容每段开头留四个字符空格，10.5磅，行距20 磅，段前段后0 磅。
\item 摘要内容后下空1行，10.5磅，行距20磅，段前段后0磅。
\item 摘要内容后为“Key words”，10.5磅，行距20磅，段前段后0 磅。“Key words”加粗。
\end{enumerate}
\end{enumerate}

\subsection{目录}

“目录”两字间空1个汉字符间隙，三号字，黑体，加粗，居中，，单倍行距，段前24 磅，段后18磅。

目录下方为章、节、条、款及其起始页码，章、节、条、款层次代号可以用阿拉伯数字或汉字。目录中一般只出现三级标题，一级标题为宋体，四号字，行距20 磅，段前6 磅，段后0 磅，两端对齐，页码右对齐，二级、三级标题为宋体，小四号字，行距20 磅，段前6磅，段后0 磅，两端对齐，页码右对齐，左缩进2个汉字符，序号与题名间空1个字符宽度。


\subsection{正文}

正文标题：一级标题为三号字，黑体，加粗，居中，单倍行距，段前24 磅，段后18 磅；二级标题为四号字，黑体，顶左，单倍行距，段前24 磅，段后6 磅；三级标题为小四号字，黑体，首行缩进2个汉字符，单倍行距，段前12 磅，段后6 磅。

正文：采用小四号字，宋体（英文用Times~New~Roman体，12磅），两端对齐，段落首行左缩进2个汉字符，行距20磅，段前段后0磅。

图：图名置于图的下方，五号字，宋体，居中，单倍行距，段前6磅，段后12 磅，图序与图名之间空2个汉字符。

表：表名置于表的上方，五号字，宋体，居中，单倍行距，段前6磅，段后6 磅，表序与表名之间空2个汉字符。表下方的注释为五号字，宋体，居左（英文用Times~New~Roman体10.5 磅），单倍行距。

注释：一般分为页末注（脚注）和篇末注。脚注，宋体，9磅（英文用Times~New~Roman，9磅），左对齐，单倍行距，段前段后0磅，按阿拉伯数字编号，每页须重新编号，具体格式参照参考书目格式。


\subsection{参考文献}

“参考文献”为三号字，黑体，加粗，居中，单倍行距，段前24磅，段后18 磅。

参考文献内容为五号字（英文用Times~New~Roman体，10.5磅），宋体，行距16 磅，段前段后0磅。


\subsection{其他}

“附录”、“索引”两字间空2个汉字间隙，三号字，黑体，加粗，居中，单倍行距，段前24磅，段后18磅。附录、索引的内容为五号字（英文用Times~New~Roman体，10.5 磅），宋体，行距16磅，段前段后0磅。

“致谢”两字间空2个汉字间隙，三号字，黑体，加粗，居中，单倍行距，段前24 磅，段后18磅。致谢正文部分首行缩进两字符，小四号字，仿宋，行距16 磅，段前段后0磅。

“在学期间的研究成果”为三号字，黑体，加粗，居中，单倍行距，段前24 磅，段后18 磅。在学期间的研究成果内容为五号字（英文用Times~New~Roman 体，10.5 磅），宋体，行距16 磅，段前段后0磅。
\end{document}
